%========================== REFERÊNCIAS ==========================
\newpage
\section*{Referências}
\addcontentsline{toc}{section}{REFERÊNCIAS}

\begingroup
\singlespacing
\setlength{\parindent}{0pt}
\setlength{\parskip}{0.7\baselineskip}
\raggedright

% ABNT, ordem alfabética. Modelos de cada tipo:
%
% ARTIGO EM PERIÓDICO
% SOBRENOME, N. Título do artigo. \textit{Nome do Periódico}, v. 10, n. 2,
% p. 100--115, 2024.
%
% TRABALHO EM EVENTO
% SOBRENOME, N. Título. In: NOME DO EVENTO, 15., 2024, Cidade.
% \textit{Anais} [...]. Porto Alegre: SBC, 2024. p. 9--15.
%
% LIVRO
% SOBRENOME, N. \textit{Título do livro}: subtítulo. Cidade: Editora, 2023.
%
% PÁGINA NA INTERNET
% A-FRAME. \textit{A-Frame documentation}. Disponível em:
% https://aframe.io/docs/. Acesso em: 9 nov. 2026.

% Os dez trabalhos selecionados na revisão. Autores, título, veículo,
% ano e DOI vieram das bases (data/revisao/resultados/selecionados.csv).
% FALTA em cada um: volume, número e páginas -- copiar da página do
% artigo, aberta pelo DOI.
[COMPLETAR: volume, número e páginas de cada referência abaixo]

ALVES, I.; ESTEVES, A. MeetingBenji: tackling cynophobia with virtual
reality, gamification, and biofeedback. In: ACM SYMPOSIUM ON VIRTUAL
REALITY SOFTWARE AND TECHNOLOGY (VRST), 2024. \textit{Proceedings}
[...]. 2024. DOI: 10.1145/3641825.3687747.

ANDERSSON, J.; HALLIN, J.; TINGSTRÖM, A.; KNUTSSON, J. Virtual reality
exposure therapy for fear of spiders: an open trial and feasibility
study of a new treatment for arachnophobia. \textit{Nordic Journal of
Psychiatry}, 2024. DOI: 10.1080/08039488.2023.2279643.

BINDER, F. P.; PÖHLCHEN, D.; ZWANZGER, P.; SPOORMAKER, V. I. Facing your
fear in immersive virtual reality: avoidance behavior in specific
phobia. \textit{Frontiers in Behavioral Neuroscience}, 2022.
DOI: 10.3389/fnbeh.2022.827673.

BRAMBILLA, S.; LIGABUE, M.; ABATE, S.; BOCCIGNONE, G.; BORGHESE, N. A.;
RIPAMONTI, L. A. Virtual reality game-based classification of
arachnophobia: a two-step clustering approach. In: IEEE CONFERENCE ON
GAMES (CoG), 2025. \textit{Proceedings} [...]. 2025.
DOI: 10.1109/cog64752.2025.11114142.

BRICE, D.; GIBSON, Z.; MCGUINNESS, F.; RAFFERTY, K. Using ultrasonic
haptics within an immersive spider exposure environment to provide a
multi-sensorial experience. \textit{Frontiers in Virtual Reality}, 2021.
DOI: 10.3389/frvir.2021.707731.

BUDIMAN, N. S.; KUSUMA, S.; TEJAWIJAYA, K. H.; CHOWANDA, A.; ANDERIES.
Biofeedback-augmented virtual reality as a tool for treating
arachnophobia. In: INTERNATIONAL CONFERENCE ON ICT FOR SMART SOCIETY
(ICISS), 2026. \textit{Proceedings} [...]. 2026.
DOI: 10.1109/iciss71466.2026.11715505.

BUNNELL, B. E.; SCHULER, K. R.; CRAGGS, J. G.; WELCH, B. M.; ONG, T.
Examining the efficacy of a telehealth-based virtual reality clinic in
treating adults with specific phobia: feasibility randomized controlled
trial. \textit{JMIR Formative Research}, 2026. DOI: 10.2196/84670.

KRUMNIKL, M.; NĚMEC, M.; VILÍMKOVÁ, A.; HUSER, F.; JAREMOVÁ, V. Do
virtual spiders feel real? Dissociation between subjective and
physiological fear during virtual reality exposure. \textit{Frontiers in
Virtual Reality}, 2026. DOI: 10.3389/frvir.2026.1916925.

LEMMENS, A.; AARTS, E.; DIBBETS, P. Itsy bitsy spider: fear and
avoidance (generalization) in a free-exploratory virtual reality
paradigm. \textit{Behaviour Research and Therapy}, 2024.
DOI: 10.1016/j.brat.2023.104442.

MOHAMED, H. G.; KHAN, M. N.; TAHIR, M.; ISMAT, N.; ZAFFAR, A.; NASEER,
F.; ALI, S. A predictive and adaptive virtual exposure framework for
spider fear: a multimodal VR-based behavioral intervention.
\textit{Healthcare}, 2025. DOI: 10.3390/healthcare13202617.

% ---------------------------------------------------------------
% TRABALHOS DE CONTEXTO -- ficaram fora da síntese, mas servem para
% a introdução e o referencial. Tirar o "%" SÓ do que for citado no
% texto: a lista de referências traz apenas o que foi citado.
%
% FERREIRA, T.; BATISTA, T.; LOPES, T.; MALHEIRO, P.; VITORINO, J.;
% VAZ, R.; SOUSA, J. de. Extended reality and phobia treatments: a
% systematic literature review. \textit{Procedia Computer Science},
% 2025. DOI: 10.1016/j.procs.2025.02.213.
%
% FUMERO, A.; MARRERO, R. J.; RIVERO, F.; ALVAREZ-PÉREZ, Y.;
% BETHENCOURT, J. M.; GONZÁLEZ, M.; PEÑATE, W. Neuronal correlates of
% small animal phobia in human subjects through fMRI: the role of the
% number and proximity of stimuli. \textit{Life}, 2021.
% DOI: 10.3390/life11040275.
%
% MARTINS, S.; VAIRINHOS, M. Ludic and narrative immersion in virtual
% reality exposure therapy to animal phobias: a systematic literature
% review. \textit{Virtual Worlds}, 2023.
% DOI: 10.3390/virtualworlds2040018.
% ---------------------------------------------------------------

\endgroup
